%% ============================================================================
%% Fichier  : annexes/Ann_A_glossaire.tex
%% Titre    : Glossaire
%% Version  : 1.0 -- novembre 2026
%% ============================================================================
%% Définitions abrégées des notions du cours ; le chapitre de référence est
%% indiqué. Alignée sur l'annexe A des autres cours de l'auteur.
%% ============================================================================

\chapter{Glossaire}
\label{ann:glossaire}

Ce glossaire reprend, sous une forme abrégée, les notions définies dans les
chapitres. Le renvoi indique le chapitre où la notion est présentée et
illustrée~; c'est là qu'il faut lire la définition complète.

\begin{description}
    \item[Acquisition] Production d'une copie fidèle d'un support ou de données,
          vérifiable par empreinte (ISO/IEC~27037~; Ch.~\ref{chap:acq}).
    \item[Acquisition logique, physique] Copie des fichiers visibles par le
          système, ou copie bit à bit de tout le support, espace non alloué
          compris (Ch.~\ref{chap:acq}).
    \item[Admissibilité] Qualité d'une preuve numérique qui la rend recevable
          dans une procédure judiciaire (Prologue~; Ch.~\ref{chap:custody}).
    \item[AFRIPOL] Mécanisme de coopération policière de l'Union africaine
          (Ch.~\ref{chap:norm-glob}).
    \item[Analyse] Examen des données acquises pour établir des faits,
          séparé de leur interprétation (Ch.~\ref{chap:PIU}).
    \item[Analyse des hypothèses concurrentes (ACH)] Méthode qui confronte
          chaque élément de preuve à toutes les hypothèses, pour réfuter
          plutôt que confirmer (Heuer~; Ch.~\ref{chap:PIU}).
    \item[Biais de confirmation] Tendance à retenir les éléments qui
          confirment une hypothèse déjà formée (Ch.~\ref{chap:deonto}).
    \item[Bloqueur d'écriture (\textit{write blocker})] Dispositif matériel ou
          logiciel qui permet de lire un support en interdisant toute écriture
          (Ch.~\ref{chap:acq}).
    \item[Carving] Récupération de fichiers à partir de leurs signatures,
          sans s'appuyer sur le système de fichiers (Ch.~\ref{chap:acq}).
    \item[Chaîne de custody (chaîne de conservation)] Registre chronologique
          de chaque personne ayant détenu une pièce et de chaque action
          effectuée, de la collecte à l'audience (Prologue~;
          Ch.~\ref{chap:custody}).
    \item[CLOUD Act] Loi des États-Unis (2018) qui permet aux autorités
          américaines d'exiger d'un fournisseur des données qu'il contrôle,
          où qu'elles soient stockées (Ch.~\ref{chap:leg-mond}).
    \item[Constatation] Fait vérifiable par un autre expert, avec sa source~;
          distincte de l'analyse (Ch.~\ref{chap:PIU}).
    \item[Convention de Budapest] Convention du Conseil de l'Europe sur la
          cybercriminalité (2001), premier traité international en la matière
          (Ch.~\ref{chap:norm-glob}).
    \item[Convention de Malabo] Convention de l'Union africaine sur la
          cybersécurité et la protection des données à caractère personnel
          (2014) (Ch.~\ref{chap:norm-glob}).
    \item[Copie de travail] Copie de l'image de référence sur laquelle porte
          l'analyse~; l'original reste scellé (Ch.~\ref{chap:custody}).
    \item[DFRWS] \textit{Digital Forensic Research Workshop}~; son rapport de
          2001 propose le premier modèle de processus de la discipline
          (Ch.~\ref{chap:meth-std}).
    \item[Données volatiles] Données qui disparaissent à l'extinction de
          l'appareil~: mémoire vive, connexions, processus (Ch.~\ref{chap:acq}).
    \item[Empreinte (\textit{hash})] Valeur de taille fixe calculée par une
          fonction de hachage cryptographique~; toute modification des
          données la change (Ch.~\ref{chap:fond-theo}, \ref{chap:custody}).
    \item[Entraide judiciaire (MLAT)] Procédure par laquelle un État demande à
          un autre d'accomplir un acte d'enquête (Ch.~\ref{chap:norm-glob}).
    \item[Format E01] Format d'image forensique qui embarque métadonnées,
          compression et empreintes (Ch.~\ref{chap:acq}).
    \item[Format brut (\textit{raw}, dd)] Copie bit à bit sans métadonnées
          (Ch.~\ref{chap:acq}).
    \item[Identification] Première étape~: reconnaître l'incident et en
          définir le périmètre (Ch.~\ref{chap:meth-std}).
    \item[Image forensique] Copie bit à bit d'un support, avec ses empreintes
          (Ch.~\ref{chap:acq}).
    \item[Intégrité] Garantie que la preuve n'a pas été modifiée depuis sa
          collecte (Prologue~; Ch.~\ref{chap:custody}).
    \item[INTERPOL] Organisation internationale de police criminelle~; ses
          bureaux centraux nationaux relaient les demandes
          (Ch.~\ref{chap:norm-glob}).
    \item[Investigation numérique] Discipline scientifique appliquée qui
          identifie, préserve, analyse et présente des preuves numériques de
          façon recevable (Prologue~; Ch.~\ref{chap:phi-fond}).
    \item[Métadonnées] Données qui décrivent d'autres données~: dates, auteur,
          logiciel, géolocalisation (Ch.~\ref{chap:gr-affaires}).
    \item[Méthode PEACE] Méthode d'entretien d'enquête en cinq étapes
          (préparation, engagement, récit, clôture, évaluation), fondée sur
          la recherche d'informations plutôt que sur l'aveu
          (Ch.~\ref{chap:deonto}).
    \item[Modèle de Dolev-Yao] Modèle d'adversaire qui contrôle entièrement le
          réseau sans pouvoir casser les primitives cryptographiques
          (Ch.~\ref{chap:fond-theo}).
    \item[Normes ISO/IEC 27037, 27041, 27042, 27043] Lignes directrices sur la
          preuve numérique~: collecte et préservation, validation des méthodes,
          analyse et interprétation, processus d'investigation
          (Ch.~\ref{chap:meth-std}).
    \item[Opposabilité] Caractère d'une preuve qui peut être valablement
          opposée à une partie en justice (Ch.~\ref{chap:phi-fond}).
    \item[Ordre de volatilité] Règle qui impose d'acquérir d'abord les données
          les plus volatiles (RFC~3227~; Ch.~\ref{chap:acq}).
    \item[OSINT] Collecte et analyse d'informations issues de sources
          légalement accessibles au public (Ch.~\ref{chap:etat-art}).
    \item[Préservation] Protection de la preuve contre toute altération, et
          traçabilité de chaque accès (Ch.~\ref{chap:custody}).
    \item[Principe de Locard numérique] Tout processus informatique laisse des
          traces dans son environnement d'exécution
          (Ch.~\ref{chap:fond-theo}).
    \item[Principes de l'ACPO] Quatre principes britanniques~: ne pas modifier
          les données, compétence de qui y accède, piste d'audit,
          responsabilité de l'enquêteur (Ch.~\ref{chap:meth-std}).
    \item[Processus investigatif universel (PIU)] Méthode du cours en phases
          séparées par des contrôles de qualité (\textit{Quality Gates})
          (Ch.~\ref{chap:PIU}).
    \item[Régime de vérité numérique] Ensemble historiquement situé des
          traces tenues pour preuves et des procédures qui les valident
          (Ch.~\ref{chap:histoire}).
    \item[Reproductibilité] Propriété d'une méthode qui, appliquée par un
          autre investigateur qualifié aux mêmes données, donne le même
          résultat (Prologue~; Ch.~\ref{chap:meth-std}).
    \item[RGPD] Règlement (UE) 2016/679 sur la protection des données
          personnelles (Ch.~\ref{chap:leg-mond}).
    \item[Scellé] Conditionnement qui garantit qu'une pièce n'a pas été
          ouverte depuis sa fermeture (Ch.~\ref{chap:custody}).
    \item[Scène de crime numérique] Tout environnement numérique susceptible
          de contenir des traces probantes d'un acte délictueux (Prologue).
    \item[SIM swapping] Fraude qui consiste à faire transférer le numéro d'une
          victime sur une nouvelle carte SIM (Ch.~\ref{chap:etat-art}).
    \item[Trace numérique] Marque laissée sans intention par un processus
          informatique, que l'investigateur interprète (Ch.~\ref{chap:phi-fond}).
    \item[Trilemme CRO] Tension formalisée entre confidentialité, fiabilité et
          opposabilité d'une preuve (Ch.~\ref{chap:phi-fond}~;
          \cite{minka2025cro}).
    \item[Volatility] Cadriciel d'analyse de la mémoire vive
          (Ch.~\ref{chap:acq}).
\end{description}
